\section{问题重述}

\subsection{问题背景}

大语言模型的训练可以用几个相互制约的资源量来刻画：参数量 $N$、训练数据量 $D$
（以 Token 计）、训练算力 $C$（以浮点运算次数 FLOPs 计）、数据质量 $Q$ 与领域配比
$\boldsymbol{p}$。标度律刻画验证集交叉熵损失 $L$ 随规模变化的经验关系，常用形式为
\begin{equation}
  L(N,D)=E+A N^{-\alpha}+B D^{-\beta},
  \label{eq:classic}
\end{equation}
其中 $E$ 为不可约损失，$A,\alpha,B,\beta$ 为待估参数\cite{hoffmann2022chinchilla}。
高质量语料日益稀缺，而扩大规模、提升数据质量、延长上下文长度都要消耗算力。
赛题要求利用真实公开数据与已标注的半合成补充数据，建立规模、质量、配比与损失之间的
定量关系，研究算力约束下的资源配置，并用历史评测数据检验规律、预测能力前沿。

\subsection{问题提出}

四个问题构成一条递进主线，前一问的输出是后一问的输入。

\textbf{问题一：数据质量评价、质量冲突消解与领域配比建模。}利用质量信号数据
（A1--A3）对 22 个质量指标进行预处理与方向统一，建立综合质量评分 $Q$，给出域级评分并
对照抽样集与扩展集；定义质量冲突、分析其成因并在扩展集上检验；利用配方实验数据
（A4--A15）建立 17 域配比 $\boldsymbol{p}$ 与交叉熵损失的定量关系，在检验集 A6--A11
上验证，并用外推表 A12--A15 讨论外推结论的稳健性。

\textbf{问题二：跨维度数据融合与广义标度律。}以 B1 为主要拟合数据估计经典标度律
\eqref{eq:classic}，用 B2--B5 做模型族外与文献验证；在此基础上构造同时包含
$N,D,Q,\boldsymbol{p}$ 的广义标度律，分析各因素的边际效用与弹性，推导提升数据质量与
扩大模型规模可相互替代的条件。半合成数据（B6--B8）不得表述为直接实验观测。

\textbf{问题三：算力约束下的多维资源联合优化与结构性转移。}在总算力不超过预算 $C$
的条件下，考虑训练开销、数据质量提升开销与长文本注意力开销，求解 $N,D,Q,\boldsymbol{p}$
的配置；上下文长度 $L_{\mathrm{ctx}}$ 外生给定，取值依据 C7；分析预算跨越不同量级时
最优配置是否发生结构性转移。

\textbf{问题四：技术演进分析与前沿预测。}结合附件 C 分离并量化规模扩张与非规模技术进步
各自的贡献，说明综合能力度量、开源筛选口径、模型类型划分与时间轴口径；建立损失与
下游评测得分之间的映射并讨论映射误差；在算力增长放缓的情景下预测未来 12 个月或
24 个月开源大语言模型能力的前沿边界，并给出不确定性分析。

\subsection{数据可信度分级}

附件中的数据来源性质不同，本文在全文中区分四类证据并分别表述：
\begin{enumerate}
  \item \textbf{观测数据}：由真实实验或公开评测直接得到的记录，如 A 组质量信号、
        B4 与 B5、C 组排行榜与元数据；
  \item \textbf{半合成数据}：以真实数据校准的补充数据，如 B6--B8，只用于补充分析；
  \item \textbf{估算参照}：由拟合面推算而非观测的数值，如 A13、A15，不能作为独立验证；
  \item \textbf{外推}：在模型识别范围或有效域之外的评估，结论须注明外推距离与条件。
\end{enumerate}
